\section{The acute orthogonal case}\label{sec:elimination}

We prove Proposition~\ref{prop:acute}. Let $L=\hat L(u,v)$ be as there and
suppose that $D\Ctrl(L)$ is not onto. All polynomial identities below were
verified by computer algebra and are checked in Lean.
% Lean: Gluck.Discrete.NativeOppositePure.fderiv_nativeLengthController_surjective_of_acute_orthogonal

\emph{Projective coordinates.} By Lemma~\ref{lem:core}(i),
$\cos\alpha_0=v\hat L_2$, $\cos\alpha_3=u\hat L_1$, $\cos\gamma_1=v\hat L_3$ and
$\cos\gamma_2=u\hat L_0$, hence
\begin{equation}\label{eq:tan}
  \tan\alpha_0\tan\gamma_2=\tan\alpha_3\tan\gamma_1=\frac{q}{uv}>1 .
\end{equation}
% Lean: Gluck.Discrete.NativeOppositePure.coreHalfLengths_angle_data,
%       coreHalfLengths_angles_acute, coreHalfLengths_tan, coreHalfLengths_tan_products,
%       coreHalfLengths_tan_product_gt_one
Put $x=\tan\alpha_0/\tan\gamma_1$, $y=1/(\tan\gamma_1\tan\gamma_2)$,
$w=1/\tan\gamma_2$ and $m=w^2$. By~\eqref{eq:tan}, $\tan\alpha_0=xw/y$,
$\tan\alpha_3=x/w$, $\tan\gamma_1=w/y$ and $\tan\gamma_2=1/w$. Since
$T<\pi/2$, $\tan\gamma_1\tan\gamma_2<1$, so $y>1$; and
$x/y=\tan\alpha_0\tan\gamma_2>1$, so $x>y>1$. With
$\lambda_*=\cos\alpha_0\cos\alpha_3/y>0$ one finds
\begin{equation}\label{eq:cos}
  \cos S=\lambda_*\,(y-x^2),\qquad \frac{\cos T}{q}=\lambda_*\,x\,(y-1);
\end{equation}
the first is $\cos S=\cos\alpha_0\cos\alpha_3(1-\tan\alpha_0\tan\alpha_3)$, and
the second uses $q\sin\alpha_3\cos\alpha_0=\cos\gamma_1\sin\gamma_2$, which
follows from the cosine relations above.
% Lean: Gluck.Discrete.NativeOppositePure.acute_projective_coordinates,
%       coreHalfLengths_normalized_cos

\emph{The equations.} By Lemma~\ref{lem:rank}, $AD=BC$. Put $a=A/2$, $b=C/2$
and $\tau=A/B$. Then $B=2a/\tau$ and $D=2b/\tau$, and since $0<p,l<1$ and
$t,h<q$, all of $A,B,C,D$ lie in $(0,1)$:
\begin{equation}\label{eq:ranges}
  0<a<\tfrac12,\quad 0<b<\tfrac12,\quad 2a<\tau,\quad 2b<\tau .
\end{equation}
% Lean: Gluck.Discrete.NativeOppositePure.incident_projective_parameters
Substituting~\eqref{eq:cos} into the second equation of~\eqref{eq:minors}
and dividing by $\lambda_*B$ gives
\begin{equation}\label{eq:affine}
  \mathcal A:=-(x+1)(x-y)(1+\tau)+2(x^2-y)b-2x(y-1)a=0 .
\end{equation}
% Lean: Gluck.Discrete.NativeOppositePure.projective_affine_of_incident_minor
Since $\tan\arcsin(c\sin z)=\varsigma_c(z)\tan z$, the four mixed arcsines
have tangents $(1-2a)\,xw/y$, $(1-2b)\,w/y$, $(1-2b/\tau)/w$ and
$(1-2a/\tau)\,x/w$. The equations $\varepsilon_1=0$ and $\varepsilon_3=0$ say
that four angles with nonzero cosines sum to $\pi$, and for such angles
$(\tan z_1\tan z_2-1)(\tan z_3+\tan z_4)=(\tan z_1+\tan z_2)(1-\tan z_3\tan z_4)$.
Clearing denominators gives
\begin{align}
  \mathcal E_1&:=y(x-m)\bigl(x(1-2a)+(1-2b)\bigr)
     -(x+1)\bigl(y^2-xm(1-2a)(1-2b)\bigr)=0,\label{eq:E1}\\
  \mathcal E_2&:=(xm-y^2)\,\tau\bigl(\tau-2b+x(\tau-2a)\bigr)
     -y(x+1)\bigl(m\tau^2-x(\tau-2b)(\tau-2a)\bigr)=0.\label{eq:E2}
\end{align}
% Lean: Gluck.Discrete.NativeOppositePure.tan_angleAtom_eq_slope_mul_tan,
%   tangent_turn_polynomial, first_projective_turn_polynomial,
%   second_projective_turn_polynomial, incidentProjectiveRankLoss_of_angles,
%   IncidentProjectiveRankLoss  (Lean names: P = x, Q = y, B = m, T = tau, r = a, k = b)
% Lean: Gluck.Discrete.NativeIncidentProjective.turnOne, turnTwo, affineRelation
So Proposition~\ref{prop:acute} follows from the next lemma.

\begin{lemma}\label{lem:elim}
There are no reals $x,y,m,a,b,\tau$ with $1<y<x$, $m>0$,
$a>0$, $0<b<\tfrac12$, $2a<\tau$ and $2b<\tau$ satisfying
\eqref{eq:affine}, \eqref{eq:E1} and~\eqref{eq:E2}.
\end{lemma}
% Lean: Gluck.Discrete.NativeIncidentProjective.not_projective_turn_solution
% Lean: Gluck.Discrete.NativeOppositePure.IncidentProjectiveRankLoss.false

\begin{proof}
Put
\[
  D_0=x(x+1)a-(x^2-y),\qquad N_0=(x+1)(x-y)+2x(y-1)a,\qquad
  \mathcal H=2bD_0+N_0 .
\]
\emph{Step 1: $\mathcal H=0$.} Write $\mathcal E_1=c_1m+d_1$ and
$\mathcal E_2=c_2m+d_2$; then $c_1d_2-d_1c_2=0$. Here
$c_1=-y(x(1-2a)+1-2b)+x(x+1)(1-2a)(1-2b)$,
$d_1=xy(x(1-2a)+1-2b)-(x+1)y^2$,
$c_2=x\tau(\tau-2b+x(\tau-2a))-y(x+1)\tau^2$ and
$d_2=-y^2\tau(\tau-2b+x(\tau-2a))+xy(x+1)(\tau-2b)(\tau-2a)$. Solving
\eqref{eq:affine} for $\tau$ (note $(x+1)(x-y)\ne0$) and substituting, one has
the identity of rational functions
\begin{equation}\label{eq:resultant}
  (x+1)(x-y)\,(c_1d_2-d_1c_2)\big|_{\tau=\tau(a,b)}
  =8xy\,(a-b)\,\bigl(b(x^2-y)+x^2a(y-1)\bigr)\,\mathcal H.
\end{equation}
% Lean: Gluck.Discrete.NativeIncidentProjective.elimination_identity (private),
%   projectiveH_eq_zero, projectiveH, projectiveD0, projectiveN
If $a=b$, then \eqref{eq:affine} reads $(x+1)(x-y)(2a-1-\tau)=0$, so
$\tau=2a-1<2a$, contradicting $\tau>2a$. Also $b(x^2-y)+x^2a(y-1)>0$. Hence
$\mathcal H=0$.

\emph{Step 2: two consequences.} The identities
\[
  (x+1)\bigl((x-y)\tau-2bxa\bigr)+\mathcal H+\mathcal A=0,
  \qquad
  D_0\,\mathcal E_1+(x+1)\,\mathcal C=c_H\,\mathcal H,
\]
with $c_H=2mx^2a-mx^2+2mxa-mx+my-xy$ and
$\mathcal C=m\,x\,\mathcal D(a)+y\,\mathcal P(a)$, where
\begin{align*}
  \mathcal P(a)&=2x^3a^2+(-3x^3+x^2y+x^2+xy)\,a+(x^3-x^2y-x^2+y^2),\\
  \mathcal D(a)&=2x(x+y-1)\,a^2+(-x^2-xy+3x-y)\,a+(y-x),
\end{align*}
give $(x-y)\tau=2bxa$ and $\mathcal C=0$.
% Lean: Gluck.Discrete.NativeIncidentProjective.projective_T_relation,
%   projective_polyC_eq_zero, polyA (= calP), polyD (= calD), polyC (= calC)

\emph{Step 3: a forced interval.} $N_0>0$, and $2bD_0=-N_0$ gives $D_0<0$.
From $(x-y)\tau=2bxa$ and $\tau>2b$ we get $a>(x-y)/x$. From $\mathcal H=0$,
$N_0+D_0=(1-2b)D_0<0$, and $N_0+D_0=x\bigl(a(x+2y-1)-(y-1)\bigr)$. Thus
\begin{equation}\label{eq:interval}
  a_-:=\frac{x-y}{x}<a<a_+:=\frac{y-1}{x+2y-1},
\end{equation}
% Lean: Gluck.Discrete.NativeIncidentProjective.projective_forced_interval
and in particular $a_-<a_+$, which is equivalent to $x^2-2y^2+y<0$.

\emph{Step 4: a sign obstruction.} $\mathcal P$ and $\mathcal D$ are convex
quadratics in $a$, so it suffices that each is negative at two points
bracketing~\eqref{eq:interval}. Direct evaluation gives
\begin{align*}
  \mathcal P(a_-)&=-xy(x-y), &
  \mathcal P\Bigl(\tfrac{x^2-y}{x(x+1)}\Bigr)&=-\frac{x(x+y)(x^2-2y+1)}{(x+1)^2},\\
  \mathcal D(a_-)&=\frac{(x-y)(x^2-xy-2y^2+y)}{x}, &
  \mathcal D(a_+)&=-\frac{y(x+y)(x^2-2y+1)}{(x+2y-1)^2}.
\end{align*}
Here $x^2-2y+1>2x-2y>0$, $x^2-xy-2y^2+y<x^2-2y^2+y<0$, and
$a_+<\frac{x^2-y}{x(x+1)}$, which after cross-multiplying is
$(x+y)(x^2-2y+1)>0$. So $\mathcal P(a)<0$ and $\mathcal D(a)<0$, hence
$\mathcal C<0$, contradicting $\mathcal C=0$.
\end{proof}
% Lean: Gluck.Discrete.NativeIncidentProjective.polyA_at_left, polyA_at_outer_right,
%   polyD_at_left, polyD_at_right, quadratic_neg_between (private),
%   polyA_neg_and_polyD_neg_of_interval, polyC_neg_of_interval
